\documentclass{article}
\usepackage[T1]{fontenc}
\usepackage{lmodern}
\usepackage{graphicx}
\usepackage{amsmath,amssymb}
\usepackage[a4paper,margin=2cm]{geometry}
\usepackage[absolute,overlay]{textpos}
\setlength{\TPHorizModule}{1mm}
\setlength{\TPVertModule}{1mm}
\usepackage[indonesian]{babel}
\usepackage{hyperref}
\setlength{\emergencystretch}{2em}
% unit: Halaman 49

\begin{document}

% === Halaman 49 (python/native) ===

Contoh: Misalkan y2  x3 – x  + 188  (mod 751) . Kurva eliptik ini memiliki N = 727 

buah titik. 

• Misalkan karakter yang akan dikodekan adalah huruf ‘B’, yang dikodekan menjadi 

nilai 11. 

• Pilih k = 20, maka x = mk + 1 = (11)(20) + 1 = 221. Sulihkan x = 221 ke dalam kurva 

eliptik y2  x3 – x  + 188  (mod 751)  456 (mod 751). Tidak ada nilai y yang 

memenuhi. 

• Coba untuk x = mk + 2 = (11)(20) + 2 = 222. Sulihkan x = 222 ke dalam kurva eliptik 

y2  x3 – x  + 188  (mod 751) . Juga tidak ada nilai y yang memenuhi. 

• Coba untuk untuk x = mk + 3 = (11)(20) + 3 = 223. Sulihkan x = 223 ke dalam kurva 

eliptik y2  x3 – x  + 188  (mod 751) . Juga tidak ada nilai y yang memenuhi. 

• Coba untuk untuk x = mk + 4 = (11)(20) + 4 = 224. Sulihkan x = 224 ke dalam kurva 

eliptik y2  x3 – x  + 188  (mod 751). Diperoleh y =  248. Jadi, karakter ‘B’ dikodekan 

menjadi titik (224, 248) pada kurva eliptik. 

• Pada proses decoding, hitung m = (x – 1)/k = (224 – 1)/20 = 11.15 = 11. Jadi 

pesan semula adalah huruf ‘B’.

49

\end{document}
